% =====================================================================
%  paper_outline_v1.tex  --  short outline of the paper (2-3 pages + toc)
%  One short description per paragraph. Long version with numbers, tables
%  and TODOs: paper_summary_v1.tex. Sources: paper/factsheets/.
% =====================================================================
\documentclass[11pt,a4paper]{article}

\usepackage[T1]{fontenc}
\usepackage{lmodern}
\usepackage[a4paper,top=2cm,bottom=2cm,left=2cm,right=2cm]{geometry}
\usepackage{amsmath,amssymb}
\usepackage{booktabs}
\usepackage{xcolor}
\usepackage{enumitem}
\usepackage{hyperref}

\setlist{itemsep=1pt,topsep=2pt}
\newcommand{\from}[2]{#1$\leftarrow$#2}
\newcommand{\open}[1]{\textcolor{red}{[#1]}}

\title{Contrastive self-supervision for jet flavour tagging:\\
augmentation strength, intersample attention and training strategies\\[4pt]
\large Paper outline v1}
\author{Riccardo Riva}
\date{November 2026}

\begin{document}
\maketitle
\tableofcontents
\newpage

\noindent Two studies. \textbf{Study I}: multi-task training (CE + contrastive loss on the jet
embedding, JC) as a function of the augmentation strength $\lambda$, for our model and for
GN2. \textbf{Study II}: contrastive pretraining (P0, P1) followed by training from scratch
(SC), fine-tuning (FT) or frozen encoders (FR). Red brackets: open points.

% ---------------------------------------------------------------------
\section{Introduction}
\begin{enumerate}
  \item \textbf{Flavour tagging.} What $b$/$c$-tagging is, why it matters at the LHC (Higgs,
        top, searches), and the physics it exploits (long-lived hadrons, displaced tracks).
  \item \textbf{Transformer taggers and their embedding.} From DL1 to GN1/GN2 and ParT/UParT;
        all map the jet to a latent embedding, shaped only indirectly by the classification loss.
  \item \textbf{Jets as tabular data.} Jet and track features as a table; tabular models
        (SAINT) embed each feature separately and add attention across features and across samples.
  \item \textbf{Contrastive self-supervision.} InfoNCE between augmented views shapes the
        embedding directly; in jet physics so far mostly with physics-motivated augmentations.
  \item \textbf{This work and contributions.} The two questions (multi-task vs pretraining),
        the contributions list, and the outline of the paper.
\end{enumerate}

% ---------------------------------------------------------------------
\section{Related work}
\begin{enumerate}
  \item \textbf{Transformer flavour tagging.} GN1/GN2 (Salt), DeepJet, ParticleNet, ParT,
        UParT; none shapes its embedding with SSL or lets a jet attend to other jets.
  \item \textbf{Deep learning on tabular data.} Trees vs deep models, TabTransformer,
        FT-Transformer, SAINT, non-parametric transformers, Set Transformer, tabular FMs.
  \item \textbf{Contrastive learning.} CPC/SimCLR/SupCon, alignment and uniformity, InfoMin
        (role of the views), tabular augmentations (CutMix, Mixup, SCARF), dimensional collapse.
  \item \textbf{SSL in jet physics.} JetCLR, RS3L, masked particle modelling, OmniJet, OmniLearn,
        contrastive approaches to mismodelling; how our generic augmentations differ.
\end{enumerate}

% ---------------------------------------------------------------------
\section{Dataset}
\begin{enumerate}
  \item \textbf{JetSet sample.} ATLAS Open Data, simulated $t\bar t$ at 13.6\,TeV, about 192M
        jets after $20<p_\mathrm{T}<250$\,GeV, $|\eta|<2.5$; truth flavour labels $b$, $c$, light, $\tau$.
  \item \textbf{Preprocessing and splits.} UPP resampling to the $c$-jet spectrum, fixed flavour
        fractions; dataset A (52M train) for study I, dataset B (40M pretrain, 18.6M train) for study II.
  \item \textbf{Inputs.} 2 jet features and 19 track features (10 continuous, 9 categorical hit
        counts), the same as GN2; about 7 tracks per jet.
\end{enumerate}

% ---------------------------------------------------------------------
\section{Model}
\begin{enumerate}
  \item \textbf{Per-feature tokens.} Every feature gets its own embedding (MLP for continuous,
        lookup table for categorical), unlike GN-like taggers that embed the whole track at once.
  \item \textbf{Column and intersample attention.} Attention across the features of an object,
        then across objects of the batch (jets with jets, tracks with tracks of all jets).
  \item \textbf{Aggregator and head.} GN2-like transformer over the tracks of each jet, pooled and
        joined with the jet features into the global jet embedding; MLP head with 4 classes.
  \item \textbf{Two variants.} Model A (ISAB row attention, 2.2M parameters, study I) and Model B
        (chunked row attention with QK-norm, 7.5M parameters, study II). \open{one model or two?}
  \item \textbf{GN2-like reference.} GN2 trained in Salt on the same jets, classification only;
        the production GN2v01 scores are used only as an external yardstick.
\end{enumerate}

% ---------------------------------------------------------------------
\section{Training objectives and strategies}
\begin{enumerate}
  \item \textbf{Augmented views.} A clean view and a view corrupted by CutMix + Mixup with
        strength $\lambda$ (the ``noise parameter'').
  \item \textbf{Contrastive losses.} Symmetric InfoNCE between the two views: on the global jet
        embedding (JC) and on the per-track and jet-feature embeddings (pretraining only).
  \item \textbf{Multi-task training (study I).} $\mathcal{L}=\mathcal{L}_\mathrm{CE}+0.3\,\mathcal{L}_\mathrm{JC}$
        from scratch, CE-only baseline, scan over $\lambda$.
  \item \textbf{Pretraining recipes (study II).} P0: contrastive on track and jet-feature
        embeddings (aggregator untrained); P1: P0 + JC (aggregator trained). 40M jets, no labels.
  \item \textbf{Downstream strategies.} SC, FT (pretrained parts at lr/10), FR (encoders frozen),
        each with or without JC. Table~\ref{tab:kinds} lists the kinds of training done.
  \item \textbf{Metrics.} Embedding quality (uniformity, effective rank); $b$-tagging
        discriminant, light/$c$/$\tau$ rejections with optimised $f_c$, $f_\tau$; seeds and
        selection rules.
\end{enumerate}

\begin{table}[h]
\centering\small
\caption{Kinds of training done with P0 and P1 (study II).}
\label{tab:kinds}
\begin{tabular}{l|cc|cc|cc}
\toprule
 & \multicolumn{2}{c|}{From scratch} & \multicolumn{2}{c|}{Fine-tuning} & \multicolumn{2}{c}{Frozen encoders} \\
Pretraining & CE & CE+JC & CE & CE+JC & CE & CE+JC \\
\midrule
none & \checkmark & \checkmark & -- & -- & -- & -- \\
P0   & -- & -- & \checkmark & -- & \checkmark & -- \\
P1   & -- & -- & \checkmark & \checkmark & \checkmark & \checkmark \\
\bottomrule
\end{tabular}
\end{table}

% ---------------------------------------------------------------------
\section{Results I: multi-task training and augmentation strength}
\begin{enumerate}
  \item \textbf{$\lambda$ scan, our model.} Weak augmentation hurts; strong augmentation
        ($\lambda\gtrsim0.55$) delays overfitting and gives the best validation CE.
        \open{batch/schedule confound, dropout baseline}
  \item \textbf{Same scan for GN2.} The contrastive term is neutral for GN2 at every $\lambda$,
        also in the HPO studies; it delays the best epoch but does not improve it.
  \item \textbf{Embedding quality.} JC drives uniformity to its floor and raises the effective
        rank; among JC runs a lower rank goes with better CE: metrics are diagnostics, not targets.
  \item \textbf{Performance vs GN2.} With optimised fractions the best $\lambda$ gives about
        +2\% $c$-rejection over GN2 at equal light and $\tau$ rejection; CE only is below GN2.
\end{enumerate}

% ---------------------------------------------------------------------
\section{Results II: pretraining and training strategies}
\begin{enumerate}
  \item \textbf{Pretraining.} P0 and P1 train stably with high-rank, uniform embeddings;
        JC-only pretraining (P2) collapses the jet-feature embedding.
  \item \textbf{Training strategies (best validation CE).} Fine-tuning gives the best models:
        \from{FT}{P0} with dropout is the best overall, then \from{FT}{P1}, both below training
        from scratch; fine-tuning always beats frozen encoders. Without dropout P1 beats P0
        (FT and FR); with dropout only P0 was run. The test rejections follow the same ranking
        on average. \open{\from{FT}{P1} with dropout not run}
  \item \textbf{Dropout and downstream JC.} Dropout helps only \from{FT}{P0} (it hurts scratch
        and frozen encoders); downstream JC is neutral from scratch, worse after P1, and breaks
        down with dropout. \open{dropout runs confounded: 80 vs 40 epochs, untuned lr; WSD redo}
  \item \textbf{Fewer labels.} With 1--100\% of the labels, P0 pretraining never overtakes
        training from scratch, and its gap shrinks as labels grow. \open{P1 not scanned}
  \item \textbf{Dataset size.} The CE gap between the two studies is explained by the number of
        labelled jets, not by model or batch size.
\end{enumerate}

% ---------------------------------------------------------------------
\section{Training stability}
\begin{enumerate}
  \item \textbf{Attention-logit growth.} Contrastive pretraining makes the intersample attention
        logits explode (one-hot attention, gradient spikes); QK-normalisation fixes it.
  \item \textbf{Collapse of CE+JC from scratch.} Without tight gradient clipping (1.0) most runs
        collapse to the class priors; in GN2, strong JC diverges without clipping.
\end{enumerate}

% ---------------------------------------------------------------------
\section{Discussion}
\begin{enumerate}
  \item \textbf{JC as a regulariser.} The contrastive term mainly delays overfitting, which helps
        a model that overfits strongly (ours) and not GN2.
  \item \textbf{What makes pretraining work.} Only pretraining that includes the aggregator,
        followed by fine-tuning, helps; possible reasons and the test to confirm them.
  \item \textbf{InfoMin view.} $\lambda$ controls the information shared by the views
        monotonically; a proper linear-probe test is left for future work.
  \item \textbf{Limitations.} Simulation only, one process, resampled pretraining data,
        GN2 without auxiliary tasks, predictions that depend on the batch-mates.
\end{enumerate}

% ---------------------------------------------------------------------
\section{Conclusions and outlook}
\begin{enumerate}
  \item \textbf{Conclusions.} One short paragraph for each study, with the main number.
  \item \textbf{Outlook.} Robustness to data/MC and MC/MC differences, physics-motivated
        augmentations, pretraining for GN2 in Salt, larger pretraining samples.
\end{enumerate}

% ---------------------------------------------------------------------
\section*{Appendices}
\begin{enumerate}
  \item Hyper-parameters and compute.
  \item Pilot studies: architecture choice, pretraining temperature and $\lambda$, learning rates.
  \item P2 (JC-only pretraining) and the fraction-optimisation details.
\end{enumerate}

\end{document}
